\section{The rank-nine self-pair filter}\label{sec:ranknine}
\subsection{The source invariant is computed on the unchanged source}
Retain the original source over $F=\F_3((t))$,
\[
 \rho=A_1\otimes A_2,
\]
where $A_h$ is the rank-three Heisenberg representation attached to the projective equation $x_h^9+x_h=t^{-h}$, for $h=1,2$. The independent quotient and wild-irreducibility hypotheses are those of the original source fixture, not newly imposed conditions on a replacement source.

After the recorded unramified splitting extension, take a nonzero constant $v$ satisfying $v^9=-v$. Then
\[
 y_v=v^3x_h^3+vx_h,
 \qquad y_v^3-y_v=-v t^{-h}.
 \tag{23.1}\label{eq:aschar}
\]
The cancellation follows by substituting $v^9=-v$ and $x_h^9+x_h=t^{-h}$. As $v$ varies, these give the eight distinct nontrivial characters of the nine-element translation quotient. A nonzero polynomial of degree three cannot vanish on all nine translations, so the corresponding additive functional is nonzero. Distinct $v$ give distinct functionals by the same argument. Since each pole order $h$ is prime to three, the inherited Artin--Schreier ramification calculation gives break $h$.

\begin{proposition}[Self-pair conductor of the retained source]\label{prop:end-source}
Under the original independence and wild-irreducibility hypotheses,
\[
 \Sw(\End\rho)=152,
 \qquad a(\End\rho)=232.
 \tag{23.2}\label{eq:source232}
\]
\end{proposition}
\begin{proof}
Conjugation on the nine Weyl operators for $A_h$ gives each character of its projective quotient once. Hence $\End A_h$ is the sum of the trivial character and eight characters of break $h$ after the unramified extension. In $\End A_1\otimes\End A_2$, there is one trivial character, eight characters of break one, and $9\cdot8=72$ characters of break two. Independence retains the distinct quotient characters. Products of characters of different breaks have the larger break. The Swan conductor is therefore $8+2\cdot72=152$.

Unramified restriction preserves the Swan conductor. By wild irreducibility, the inertia-fixed endomorphisms are exactly the scalars. Thus $\dim(\End\rho)^{I_F}=1$ over $F$ itself. Equation~\eqref{eq:artin} gives $81-1+152=232$. No assumption that Frobenius fixes each individual quotient character is needed.
\end{proof}

\source{[Q13] supplies the source and the ramification inputs. The displayed Artin--Schreier cancellation, character separation and conductor addition are expanded here. The new exact suite constructs the nine-element constant subspace inside $\F_{81}$ and checks all its additive character values. This finite check does not reprove local ramification theory.}

\subsection{The actual-order certificate cannot be omitted}
The proposed rank-nine family has a complete principal order $\A$ of period nine, radical $\Pp$, and an actual final parameter field $K/F$ that is totally ramified of degree nine. An intermediate field $E$ has degree three. Write $\mathfrak Q$ for the radical of the induced order in the $E$-centralizer. The retained candidate certificate specifies leading level $24$, lower even critical level $c=2a$, and actual simple-character orders
\[
 \mathfrak J=\OO_K+\mathfrak Q^a+\Pp^{12},
 \qquad
 \mathfrak H=\OO_K+\mathfrak Q^{a+1}+\Pp^{13}.
 \tag{23.3}\label{eq:rank9orders}
\]
These equations are substantive construction inputs. The mere pair of numbers $(24,c)$ does not imply them.

The certificate also ensures that the $K$ contribution adds no unexpected graded direction. In the notation of [Q13], minimality of the leading approximation $\gamma$ and the correction $d\in\Pp^{-c}$ give
\[
 [\gamma,\OO_K]\subseteq\Pp^{-c},
 \qquad
 \OO_K\subseteq\mathfrak B+\Pp^{24-c}.
 \tag{23.4}\label{eq:compat}
\]
For $c=4$ and $c=10$, the last exponent exceeds thirteen. Through the grades used in \eqref{eq:rank9orders}, the $K$-grade therefore lies inside the three-dimensional centralizer grade. Its own dimension is one. This is an actual commutator and approximation assertion, not an integral splitting of a wildly ramified bimodule.

Let $I^+=1+\Pp$. The successive full grades have dimension nine, the intermediate grades dimension three, and the $K$-grades dimension one. It follows that
\begin{align}
 I&=\log_q[I^+:J^1]=8(a-1)+6(12-a)=2a+64,\tag{23.4a}\label{eq:rank9I}\\
 D&=\log_q[J^1:H^1]=2+6=8,
 \qquad \dim\eta=q^4.\tag{23.4b}\label{eq:rank9D}
\end{align}
Since $K/F$ is totally ramified of full degree, the finite residue factor is a character and the torsion number is one.

\subsection{A formal-degree calculation with its measure fixed}
Normalize Haar measure by $\operatorname{vol}(\GL_9(\OO_F))=1$ and central measure by $\operatorname{vol}(\OO_F^\times)=1$. Then
\[
 \operatorname{vol}(I^+)=\prod_{j=1}^{9}(q^j-1)^{-1},
 \qquad
 \operatorname{vol}(J)=(q-1)q^{-I}\operatorname{vol}(I^+).
\]
The compact-induction degree formula and the Steinberg degree formula give
\[
 d(\pi)=\frac{q^4}{9\operatorname{vol}(J)},
 \qquad
 d(\mathrm{St}_9)=\frac1{9}\prod_{j=1}^{8}(q^j-1),
\]
so that
\[
 \frac{d(\pi)}{d(\mathrm{St}_9)}
 =\frac{q^9-1}{q-1}\,q^{2a+68}.
 \tag{23.5}\label{eq:degree9}
\]
The factor nine comes from the actual noncompact extension $[K^\times J:F^\times J]$. Omitting it before comparing the measures changes the calculation. An unramified twist may be used to unitarize solely for the formal-degree comparison. That operation does not replace the prescribed noncompact scalar in the construction, and it leaves the self-pair unchanged.

The primary all-characteristic formulas [AP, Theorems 6.5 and 6.9] use the torsion number $r_t$ and an auxiliary parameter $\delta$:
\[
 \frac{d(\pi)}{d(\mathrm{St}_n)}
 =r_t\frac{q^n-1}{q^{r_t}-1}
 q^{(r_t-n+\delta)/2},
 \qquad
 a(\pi^\vee\times\pi)=\delta+n^2-r_t.
 \tag{23.6}\label{eq:AP}
\]
These are the named classical inputs here, not a conjectural formal-degree identity. With $n=9$, $r_t=1$ and \eqref{eq:degree9}, they give
\[
 \delta=4a+144=2c+144,
 \qquad a(\pi^\vee\times\pi)=4a+224=2c+224.
 \tag{23.7}\label{eq:pair9}
\]

\begin{center}\begin{tabular}{lrrrrr}\toprule
Certified profile & $c$ & $I$ & $D$ & $\delta$ & Self-pair $a$\\\midrule
First repair & 10 & 74 & 8 & 164 & 244\\
Deeper repair & 4 & 68 & 8 & 152 & 232\\
Retained source & --- & --- & --- & 152 & 232\\\bottomrule
\end{tabular}\end{center}

\begin{corollary}[A necessary selection inside this certified family]\label{cor:c4}
If a member of the actual-order family \eqref{eq:rank9orders} is to match the retained source under a correspondence preserving the ordinary self-pair conductor, then $c=4$. The critical-ten member is excluded. The critical-four member passes this invariant, but the comparison does not prove its source identification.
\end{corollary}
\begin{proof}
Equate \eqref{eq:source232} and \eqref{eq:pair9}. The equation $2c+224=232$ has the unique solution $c=4$. Equality of this one invariant supplies no character-occurrence or full pair-factor identity.
\end{proof}

\source{[Q13] provides the conditional actual-order certificate and calculation. [AP, formulas (17), (29), Theorems 6.5, 6.9 and Remark 6.7] were checked at their stated measures and field scope. This block does not newly construct the proposed corrected stratum. It preserves the requirement to prove its actual order equations before applying the numerical filter.}
